\section{Especificação das Interfaces Externas}

\subsection{Requisitos de Interface Externa}

\subsubsection{Interfaces do Usuário}
As interfaces do usuário do \ProjetoNome\ devem atender de forma coerente os dois perfis operacionais contemplados nesta especificação: Aluno e Professora. Como o produto é um aplicativo \textit{mobile} com foco no contexto de uso da academia, a interação deve privilegiar leitura rápida, navegação objetiva e resposta imediata nos fluxos mais frequentes, especialmente consulta de agenda, visualização do mapa, realização de \textit{check-in}, cancelamento de reserva e manutenção operacional.

Para o Aluno, a interface deve priorizar poucos passos, linguagem direta e clareza sobre o estado da aula, da posição escolhida e da disponibilidade real das \textit{bikes}. O mapa da sala deve refletir a organização física do ambiente e distinguir, de maneira compreensível, posições livres, ocupadas, indisponíveis ou afetadas por manutenção. Informações operacionais relevantes ao processo de reserva, como turma, data, horário, \textit{mix} atual e confirmação do \textit{check-in}, devem estar acessíveis de forma simples e contextualizada.

Para a Professora, a interface deve oferecer visão operacional das entidades que sustentam o fluxo principal do produto, incluindo salas, \textit{bikes}, posições, turmas, \textit{mixes}, grupos de alunos e manutenções. Nessa frente, a solução deve favorecer a compreensão da relação entre os elementos do domínio, evidenciar os impactos de ativação, inativação, manutenção e cancelamento e permitir acompanhamento da ocupação das aulas de forma compatível com a rotina da academia.

Independentemente do perfil, a interface não deve depender exclusivamente de cor para representar estados críticos. Situações como disponibilidade, indisponibilidade, bloqueio, ocupação e erro devem possuir reforço textual, icônico ou visual equivalente. O sistema também deve fornecer retorno explícito para ações de sucesso, falha, ausência de dados, conflito de reserva e carregamento de informações.

\subsubsection{Interfaces de Hardware}
No escopo desta especificação, o \ProjetoNome\ não depende de integração direta com hardware especializado, sensores, leitores, catracas, impressoras ou periféricos externos específicos. A interface de hardware relevante é a dos próprios dispositivos utilizados para executar o aplicativo, com destaque para \textit{smartphones} Android e dispositivos iOS utilizados pelos alunos, além de dispositivos móveis ou computadores de apoio utilizados pela Professora na operação da academia.

Os dispositivos de uso principal devem oferecer, no mínimo, tela sensível ao toque com área útil suficiente para leitura da agenda e visualização do mapa da sala, capacidade de armazenamento local para persistência dos dados do aplicativo e recursos de processamento compatíveis com renderização do mapa, consultas e operações de \textit{check-in} sem degradação perceptível do fluxo. O projeto deve ser executável em aparelhos compatíveis com as configurações mínimas de \textit{build} definidas para Android e iOS no ambiente da aplicação, sem exigir acessórios adicionais para que as funções centrais do produto sejam realizadas.

No recorte atual, não há necessidade de câmera, microfone, geolocalização, Bluetooth ou leitura de arquivos locais para execução das funções centrais. Como o produto trabalha com URLs informadas em cadastro para fotos, links e conteúdos externos, a exibição desses recursos depende apenas da capacidade normal de renderização do dispositivo e, quando aplicável, de conectividade para acesso ao endereço remoto correspondente.

\subsubsection{Interfaces de Software}
As interfaces de software do \ProjetoNome\ concentram-se principalmente na comunicação entre a camada de apresentação do aplicativo, as regras de domínio e os mecanismos de persistência local que sustentam o funcionamento do produto. No recorte atual, o sistema deve preservar consistência entre agenda, turma, mapa, posição, \textit{bike}, \textit{mix}, histórico e manutenção, refletindo nas telas exatamente o estado operacional mantido pela aplicação.

Como a solução foi concebida com foco em operação local, a principal interface de software externa ao fluxo visual corresponde ao mecanismo de armazenamento persistente utilizado pelo aplicativo. Essa interface deve permitir gravação, consulta e atualização dos dados do domínio de forma íntegra, especialmente nas operações críticas de \textit{check-in}, cancelamento e atualização da disponibilidade da turma.

No estado atual do projeto, não foram identificadas integrações obrigatórias com serviços externos de autenticação, notificações \textit{push}, captura de mídia ou sincronização remota contínua. Também não foram observadas, nos artefatos móveis principais, permissões sensíveis específicas para câmera, galeria, microfone, localização ou contatos. Assim, a operação básica do produto deve permanecer desvinculada dessas autorizações.

O produto admite, como evolução futura, interfaces complementares com serviços de autenticação, sincronização, notificações e fontes externas de informação. Entretanto, tais integrações não constituem dependências obrigatórias para o funcionamento básico especificado neste documento e não devem comprometer a autonomia dos fluxos centrais do sistema.

\subsubsection{Interfaces de Comunicação}
No recorte atual, o \ProjetoNome\ não depende de comunicação contínua em rede para executar suas funções centrais. Os fluxos principais de consulta, visualização do mapa, reserva, cancelamento e acompanhamento operacional devem permanecer utilizáveis mesmo na ausência de conectividade externa, desde que o contexto local do aplicativo esteja disponível.

As necessidades de comunicação identificadas no projeto estão associadas principalmente ao acesso eventual a endereços externos informados por URL, como imagens remotas de perfil, fotos de artistas, links musicais e referências de vídeo-aula. Esse tipo de comunicação é complementar ao fluxo principal e não deve impedir o uso do aplicativo quando indisponível, podendo apenas limitar a exibição do conteúdo remoto correspondente.

Nos artefatos móveis atualmente configurados, a permissão de internet aparece apenas nos manifestos de desenvolvimento do Android, compatível com execução e depuração durante o ciclo de implementação. No recorte funcional descrito nesta especificação, não há dependência obrigatória de comunicação em rede para reserva, agenda, mapa ou manutenção operacional.

Ainda assim, o contexto da solução admite interfaces de comunicação complementares para frentes futuras do ecossistema, como sincronização entre dispositivos, envio de notificações, integração com serviços remotos e troca de dados com componentes externos da academia. Quando presentes, essas comunicações devem ocorrer de forma controlada, sem alterar a previsibilidade do fluxo principal nem introduzir dependência obrigatória para o uso cotidiano do produto.

Sempre que houver troca de dados entre o aplicativo e componentes externos, a comunicação deverá respeitar as restrições de segurança, privacidade e finalidade aplicáveis ao domínio, preservando a integridade das informações e o tratamento adequado dos dados pessoais no contexto da LGPD.